\section{Identification and Methods}
\label{sec:methodology}

This section sets out the empirical strategy for the central question (C1): whether
the suspension changed the \emph{spot} realized volatility of the affected commodities,
relative to a credible counterfactual. We treat the exercise as a policy
\emph{evaluation} rather than as the estimation of a structural parameter: the
suspension was not randomly assigned, and we are explicit throughout about the
selection problem it creates. We deploy several estimators that fail under different,
and in places complementary, assumptions; we are careful, however, \emph{not} to rest
the identification claim on their agreement, because they share one panel and the same
five treated/five donor commodities, so that any agreement is largely mechanical rather
than independent corroboration. What credibility the point estimate has comes instead
from its stability to leaving out any single treated commodity and from the
falsification battery, both reported in Section~\ref{sec:results}. We first describe the outcome
and the panel; then the difference-in-differences (DiD) baseline; then the
synthetic-control family that supersedes it as the primary design; then the inference
machinery appropriate to a setting with very few treated clusters; then the
volatility-dynamics (GARCH) arm of the secondary question (C2); and finally the
pre-registration that disciplines these choices.

\subsection{Outcome, panel, and treatment timing}

The outcome is the log of the 30-day realized volatility of daily log spot returns,
annualized using $\sqrt{252}$ trading days, evaluated monthly: $y_{it} = \ln(\mathrm{rv30})$.
Realized volatility is computed from district-level mandi modal spot prices
(MoA Agmarknet, re-served through the CEDA Data Portal API; see Section~\ref{sec:data}),
spanning 2017-01 to 2025-10. We construct two analytical levels: a national
(district-median) series used for headline figures and per-commodity synthetic
controls, and a district panel (commodity $\times$ district $\times$ month, with on the
order of 1{,}994 commodity-district units assembled; the headline DiD is estimated on
$104{,}491$ observations across $1{,}570$ units after listwise deletion and the
within-window exclusions) used for the panel estimators and to escape the few-cluster
trap discussed below. Robustness twins replace the 30-day window with a
60-day window.

A measurement detail is load-bearing. The raw mandi series behave as seven-day
calendar grids---prices are carried across non-trading days---so that annualizing with
$\sqrt{252}$ trading days mismatches the units and injects spurious zero-return blocks.
We therefore restrict every series to weekdays before computing returns and realized
volatility. As documented in Section~\ref{sec:results}, this correction materially
changes the falsification picture; the entire analysis is run on the trading-day-clean
series.

Treatment is assigned at the \emph{commodity} level, so all inference is clustered by
commodity and never over-credits the district count. We use \emph{per-commodity
(staggered) treatment dates} taken from the policy ledger: chana on 2021-08-16,
mustard on 2021-10-08, and wheat, soybean, moong (and CPO) on 2021-12-20.\footnote{The
December 2021 SEBI directions to the exchanges were dated 19~December with the press
release on 20~December; we date the December cohort to 2021-12-20.} A commodity already
suspended may not serve as a donor for a later episode.

Two commodities receive special handling, both pre-registered. \emph{Paddy is dropped}
from the primary analysis: $40.3\%$ of its daily spot returns are exactly flat because
FCI/state MSP procurement pins the price at the support level for long stretches, so its
realized volatility is a mechanical censoring artifact rather than market volatility.
\emph{Wheat is retained but MSP-flagged} ($19.7\%$ flat returns); as a core commodity it
is kept and read as a robustness check rather than dropped. \emph{CPO has no mandi spot
series} and is therefore routed out of the spot-volatility track entirely (to the
domestic--international spread arm, H4). The primary treated set is thus
\{chana, wheat, soybean, moong, mustard\}.

The donor pool follows the pre-registered ``Option B'': a clean, screened core of
non-banned industrials and spices---\{castor, guarseed413, cotton, jeera, turmeric\}---augmented
by non-banned food/oilseed staples (barley, maize, jowar, bajra, ragi, groundnut,
sesamum, sunflower) that restore the food-inflation/MSP channel on which chana and
wheat load. Donors must lie in the same or an adjacent commodity group where possible,
remain traded and liquid through the ban, and carry no own treatment inside the
estimation window. We use the clean guar underlying \texttt{guarseed413} (correlation
$0.99$ with the futures) and exclude the gum-contaminated \texttt{guar id~75}
(correlation $0.04$); cotton's August 2022 unit-break is confined to the futures track.

\subsection{Difference-in-differences baseline}

The two-way fixed-effects baseline is
\begin{equation}
y_{it} \;=\; \alpha_i + \lambda_t + \beta\,(\text{banned}_i \times \text{post}_t) + X_{it}'\gamma + \varepsilon_{it},
\label{eq:did}
\end{equation}
with unit fixed effects $\alpha_i$, month fixed effects $\lambda_t$, and standard errors
clustered at the commodity level. The coefficient of interest is $\beta$, reported as the
proportional effect $e^{\beta}-1$. We stress that this specification is the \emph{foil},
not the headline. Its identifying assumption---parallel trends in log volatility---is
exactly what the suspension's selection-on-the-2021-run-up threatens: the literature on
selection on a transitory, mean-reverting shock \citep{Ashenfelter1978, ChabeFerret} shows
that DiD will attribute the subsequent reversion to the policy. We therefore subject
Equation~\eqref{eq:did} to the falsification battery in Section~\ref{ssec:inference} and
read it only alongside the synthetic-control family, while noting that the family shares
the same panel and donor set, so its agreement is corroboration of a weak, non-independent
kind.

\subsection{The synthetic-control family as the primary design}

Because the binding constraint is policy endogeneity together with few treated and few
clean control commodities, the primary design is the \emph{synthetic-control (SC) family},
run as several members rather than as a single estimator. The common logic is that, by
matching the treated commodity's \emph{actual} pre-ban volatility path---including the 2021
run-up that selected it for treatment---an SC estimator absorbs the very pre-trend that
biases DiD. We deploy three members of the family.

\paragraph{(i) Abadie synthetic control (commodity level).}
For each treated commodity we construct a synthetic counterfactual as a convex combination
of donor commodities, with weights $w \ge 0$, $\sum w = 1$, fitted to the pre-treatment
district-median $\ln(\mathrm{rv30})$ path \citep{Abadie2010, AbadieReview}. The treatment
effect is the post-period gap between the treated path and its synthetic. This is the
estimator most exposed to overfitting a noisy volatility target \citep{AbadieVives}, which
is why it is triangulated rather than reported alone.

\paragraph{(ii) District-unit synthetic control.}
To multiply the effective number of comparisons and exploit cross-sectional variation, we
also fit one synthetic counterfactual per treated \emph{district} unit, drawing donors from
the corresponding non-treated district series. This is the route out of the few-cluster
trap and is the design on which the in-space placebo distribution is richest; with the clean
core of five donors the placebo floor is coarse, which is the explicit motivation for
widening the donor pool with food staples.

\paragraph{(iii) Synthetic difference-in-differences.}
Synthetic DiD \citep{Arkhangelsky2021} combines unit weights (as in SC) with \emph{time}
weights that down-weight anomalous pre-periods. The time weights are the cleanest available
answer to the 2021 food-inflation run-up: they discount the very months on which the policy
was selected, while the unit weights match the cross-sectional level. We report a pooled
estimate and per-commodity estimates.

These estimators rest on different assumptions---convex-hull spanning and good pre-fit
(Abadie SC), the same plus disaggregated comparisons (district-unit SC), and a latent
factor structure with informative time weights (Synthetic DiD)---and the DiD baseline rests
on parallel trends. \citet{FermanPinto} make the honest caveat explicit: synthetic control
\emph{reduces but does not eliminate} selection bias under confounding. We therefore do not
claim that any one of these estimators is unbiased. We also do not lean on the fact that these
estimators agree on the sign and rough magnitude: because they share one panel, one outcome,
and the same five treated and five donor commodities---and because the Synthetic DiD unit
weights turn out close to uniform, so the design collapses toward the DiD---their agreement is
largely mechanical, one comparison re-weighted rather than several independent confirmations.
We report it for completeness but treat the leave-one-commodity-out stability of the point
estimate and the falsification battery, not the across-estimator agreement, as what credibility
the design carries at this cluster count.

\subsection{Inference with few treated clusters}
\label{ssec:inference}

With only five treated commodity clusters, conventional asymptotic standard errors are not
trustworthy and we layer three forms of inference.

\paragraph{Wild-cluster bootstrap.}
For the DiD coefficient in Equation~\eqref{eq:did} we report a wild-cluster (Webb six-point)
bootstrap clustered by commodity \citep{CameronGelbachMiller, MacKinnonWebb}, alongside the
analytic clustered $p$-value, and headline the bootstrap value as the more defensible quantity.

\paragraph{In-space placebos.}
For the SC family we treat each donor (or donor district) as if it were fake-treated on the
true date, recompute the gap, and locate the treated unit's effect (or pre/post RMSE ratio)
in the resulting placebo distribution \citep{Abadie2010}. With a small donor pool the attainable
placebo $p$-value floor is mechanically coarse (on the order of $1/6$ with the clean core),
which is why widening the donor pool is the decisive significance lever rather than a cosmetic step.

\paragraph{Placebo-date and pre-trend falsification.}
Two pre-registered death-tests guard against mistaking mean reversion for a policy effect.
First, a \emph{placebo-date} test re-runs the entire design with a fake suspension at
December 2019, using pre-ban data only: a design that ``finds'' an effect where none can
exist is disqualified. Second, a \emph{joint pre-trend} (event-study leads) test asks whether
the treated and counterfactual paths were already diverging before the suspension. Under the
pre-registered rules, a failed pre-trend test kills the DiD regardless of significance, and a
large placebo effect demotes the DiD magnitude in favour of the SC family. As reported in
Section~\ref{sec:results}, both tests clear their formal thresholds on the
trading-day-clean, paddy-dropped data---a reversal of their pre-fix behaviour that we
document transparently---but each only \emph{marginally}: the placebo still recovers
$56\%$ of the headline and the joint pre-trend null masks a significant most-recent
lead bin, so we read this as a marginal pass, not a clean exoneration.

\subsection{Volatility dynamics (C2): ICSS-within-regime GARCH}
\label{ssec:garch}

The secondary question (C2) concerns whether volatility \emph{dynamics}---persistence and
the level of conditional variance---shifted around the suspension. The naive design here, a
GARCH(1,1) with a ban step-dummy comparing pre- versus post-suspension persistence, is doubly
invalid in this setting. \citet{Hillebrand} and \citet{Lamoureux1990} show that an
\emph{unmodelled} variance-level break biases estimated GARCH persistence $(\alpha+\beta)$
upward toward unity, so a single level shift can masquerade as a persistence change; and the
break date is itself selected on the outcome. We therefore adopt the disciplined baseline of
\emph{ICSS-then-within-regime} estimation: we first date breaks in the unconditional variance
using the Inclán--Tiao algorithm with the HAC-corrected Sansó--Aragó--Carrión statistic
\citep{InclanTiao, Sanso}, then estimate GARCH \emph{within} regimes, and test whether any
variance break in fact sits at the suspension date rather than elsewhere (the variance-process
analogue of the C1 placebo). As reported in Section~\ref{sec:results}, this procedure finds no
variance break near the suspension on the affected series, and the district-median chana GARCH
is degenerate (persistence pinned near unity); we therefore lean on the C1 spot-volatility
estimators rather than on GARCH for the volatility verdict. The futures--spot basis arm of C2
is internally impossible as originally stated---no post-suspension futures exist for the affected
commodities---and is reframed as a domestic-versus-international spread (future work, H4), blocked
on vendor futures data.

\subsection{Pre-registration}

To prevent the specification from being reverse-engineered from the answer, the primary C1
design and the robustness curve were frozen before the synthetic-control results were
estimated (\texttt{preregistration\_c1.md}). The pre-registration locks the outcome and its
trading-day construction, the district-panel level and commodity-level clustering, the treated
set (paddy dropped, wheat MSP-flagged, CPO routed out), the per-commodity staggered dates, the
donor roster and exclusion windows from the policy ledger, the estimator family, and---crucially---the
pass/fail rules: ``ban raised volatility'' is supported only if the SC gap is positive in a
defensible band and the treated RMSE ratio lies in the top tail of the placebo distribution; a
negative gap that beats the placebos is reported as the suspension being \emph{associated with
lower} spot volatility, descriptively, without over-causalising; and the headline point estimate
is reported with its leave-one-commodity-out band and honest few-cluster inference, never as a
clean across-estimator ``convergence.'' The robustness (specification)
curve spans $\pm\{30,60\}$-day realized volatility $\times$ \{drop paddy / keep flagged\}
$\times$ \{core / core+food donors\} $\times$ \{per-commodity dates / single December 2021 date\}.

\subsection{Interpretation}

We read the resulting estimate as quasi-causal and descriptive: the suspension is
\emph{associated with} a change in spot realized volatility, with endogeneity reduced---the
placebo-date and pre-trend tests clear their thresholds, if only marginally---but not eliminated. The known limitations are stated
plainly: few treated clusters (five to six commodities); MSP censoring in the affected food
staples; the absence of any post-suspension futures for the affected commodities; and a donor
pool whose food-staple, district-level expansion is in progress and which, once complete, is
expected to sharpen both the twins and the placebo inference. (Full citations are collected in
\texttt{references.md}.)
